\section{Allocation as Evidence Acquisition}

Fusion's governing mechanism follows directly from the intelligence allocation
problem. The system does not need to predict the full demand of a trajectory
before the work begins. It can discover that demand progressively by treating
inference as a closed-loop process of evidence acquisition.

At each point in a workflow, Fusion maintains the current state of the task and
asks four questions:

\begin{enumerate}
  \item What decision-relevant uncertainty remains?
  \item What evidence could materially change the next answer or action?
  \item Which model--role or tool computation is best suited to acquire it?
  \item Is the expected value of that computation greater than its cost,
  latency, and risk?
\end{enumerate}

The resulting observation updates the task state, and the questions are asked
again. Allocation is therefore not a prelude to execution. It is a process that
continues throughout execution.

\subsection{From task assignment to computation allocation}

The distinction changes the unit of orchestration. A router assigns a request
to a model. Fusion allocates the \emph{next evidence-producing computation} in
the context of an evolving state. Its decision includes not only a model, but a
role, a bounded context view, a budget, and a criterion for stopping.

\begin{table}[H]
\centering
\caption{Routing and allocation solve different control problems.}
\small
\begin{tabular}{@{}p{0.23\textwidth}p{0.32\textwidth}p{0.36\textwidth}@{}}
\toprule
 & Model routing & Fusion allocation \\
\midrule
Unit of decision & Request & Next computation \\
Demand model & Upfront class or difficulty & Evolving information state \\
Control action & Choose an answering model & Choose source, role, view, budget, or stop \\
Output semantics & Answer & Attributed evidence or commitment \\
Time horizon & One assignment per request or turn & Closed loop within and across the trajectory \\
\bottomrule
\end{tabular}
\end{table}

This is not a claim that routing is obsolete. Routing remains a useful action
inside Fusion when a request is sufficiently bounded. The difference is that
it is no longer the governing abstraction for the whole system.

\subsection{Decision-relevant uncertainty}

Language models expose several signals associated with uncertainty. Token
entropy measures variation over expression; semantic entropy can better capture
disagreement over meaning
\citep{kuhn2023semantic,farquhar2024semantic}. Fusion must reason at a third
level: whether an unresolved variable could change the system's next
commitment.

This distinction prevents activity from masquerading as progress. Several
models can agree because they share training patterns and blind spots. Many
different answers can be irrelevant to the load-bearing decision. Conversely,
a single tool observation or counterexample can collapse an entire branch of
reasoning.

Fusion therefore values a computation by its expected effect on the decision:

\begin{equation}
  \text{value of a call}
  = \text{expected decision improvement}
  - \text{compute} - \text{latency} - \text{risk}.
\end{equation}

The system does not need a perfect Bayesian model of an open-ended task. It
needs practical, auditable approximations: unresolved claims, candidate
quality, evidence coverage, workflow readiness, failure risk, and the observed
conditional value of available sources. Each acquisition can improve these
estimates before the next allocation is made.

\subsection{Models as conditional evidence sources}

Under this view, a model is not a scalar score and a call is not a vote. A model
is a conditional source of evidence. Its contribution depends on the task,
role, context, evidence already present, provider surface, and budget. The same
model may be valuable as a critic and redundant as a second solver. A frontier
model may be indispensable for synthesis yet unnecessary for a bounded search
or check.

This is why an asymmetric model pool can outperform its apparent ingredients.
Fusion is not replacing strong intelligence with weak intelligence. It is
separating the different kinds of work hidden inside a model call and sourcing
each where it has the greatest marginal value.

\subsection{Two principles}

The evidence-acquisition view leads to Fusion's foundational design principle:

\principle{Organize computation around uncertainty, not models.}

But plural evidence creates a second problem. In an agentic workflow, multiple
models cannot independently redefine the task, mutate shared state, or issue
incompatible actions. Epistemic plurality must not become operational
fragmentation. Fusion therefore follows a second principle:

\principle{One state, many perspectives, one commit.}

Together, the principles define a compound intelligence system. The first
determines why computation is acquired. The second determines how multiple
computations can contribute without sacrificing coherence and accountability.
